\section{核岭回归与核方法的代价}

上一节的表示定理许诺：``经验损失加范数惩罚''的无穷维问题必然塌缩成 \(n\) 维问题。本节先兑现这个许诺——用于平方损失，推出核岭回归的闭式解；再把它与``正则化从 Ridge 到 Lasso''中的原始形式对照，看清同一件事的两种视角；最后清点账单：Gram 矩阵的存储与求逆代价划定了朴素核方法的规模边界，而 \(\norm{f}_{\mathcal{H}_k}\) 与泛化理论的联系则解释了为什么这笔账值得付。

\subsection{从表示定理到 n 维问题}

给定训练数据 \((x_1,y_1),\dots,(x_n,y_n)\)，核 \(k\)，正则参数 \(\lambda>0\)。核岭回归求解

\[
\min_{f\in\mathcal{H}_k}\ 
\sum_{i=1}^{n}\bigl(y_i-f(x_i)\bigr)^{2}
+\lambda n\,\norm{f}_{\mathcal{H}_k}^{2}.
\]

写成 \(\lambda n\) 是为了让闭式解干净：损失是求和而非平均，正则强度随 \(n\) 缩放。由表示定理，极小解必为 \(f=\sum_{j=1}^{n}\alpha_j k(x_j,\cdot)\)。记 Gram 矩阵 \(K=\bigl[k(x_i,x_j)\bigr]\in\R^{n\times n}\)，则

\[
f(x_i)=\sum_{j=1}^{n}K_{ij}\alpha_j=(K\alpha)_i,
\qquad
\norm{f}_{\mathcal{H}_k}^{2}=\alpha^{\trans}K\alpha,
\]

问题化为 \(n\) 维二次问题

\[
\min_{\alpha\in\R^{n}}\ 
\norm{y-K\alpha}^{2}+\lambda n\,\alpha^{\trans}K\alpha.
\]

这是 \(\alpha\) 的凸二次函数：第一项的 Hessian 是 \(2K^{2}\succeq0\)，第二项的是 \(2\lambda nK\succeq0\)。令梯度为零：

\[
-2K(y-K\alpha)+2\lambda nK\alpha=0
\iff
K\bigl[(K+\lambda nI)\,\alpha-y\bigr]=0.
\]

一个充分的取法是解括号内的方程。由于 \(K\succeq0\) 且 \(\lambda n>0\)，矩阵 \(K+\lambda nI\) 对称正定、必然可逆，于是

\[
\widehat\alpha=(K+\lambda nI)^{-1}y,
\qquad
\widehat f(x)=\sum_{i=1}^{n}\widehat\alpha_i\,k(x_i,x).
\]

（\(K\) 奇异时梯度方程还有其他解，但它们在 \(K\) 的零空间方向上彼此相差一个分量，给出的函数相同；上面的闭式解总是合法的。）训练归结为一次 \(n\times n\) 线性方程组求解，预测是 \(n\) 个核值的加权和。

\subsection{与 Ridge 原始形式的对照}

取线性核 \(k(x,z)=\ip{x}{z}\)，此时 \(K=XX^{\trans}\)，闭式解变成 \(\widehat\alpha=(XX^{\trans}+\lambda nI)^{-1}y\)，预测

\[
\widehat f(x)=\sum_{i=1}^{n}\widehat\alpha_i\,\ip{x_i}{x}
=\ip{X^{\trans}\widehat\alpha}{x}.
\]

令 \(\widehat w=X^{\trans}\widehat\alpha\)，与 Ridge 的原始解 \(\widehat w=(X^{\trans}X+\lambda nI)^{-1}X^{\trans}y\)（正则参数按同一尺度写出）对比，两者给出完全相同的预测——恒等式

\[
X^{\trans}(XX^{\trans}+\lambda nI)^{-1}
=(X^{\trans}X+\lambda nI)^{-1}X^{\trans}
\]

两端右乘 \(y\) 即见。这是同一件事的两种视角：

\begin{itemize}
\tightlist
\item
  原始形式在特征权重里思考：未知量 \(w\in\R^{p}\)，求解 \(p\times p\) 系统，适合特征不多、样本很多的情形；
\item
  对偶（核）形式在样本相似度里思考：未知量 \(\alpha\in\R^{n}\)，求解 \(n\times n\) 系统，特征维数 \(p\) 从账单中彻底消失。
\end{itemize}

多项式核、RBF 核对应的 \(p\) 巨大甚至无穷，原始形式根本写不出来，对偶形式却只用 \(n\times n\) 的 \(K\)——这正是``不构造 \(\phi\)''许诺的兑现。反过来，样本量远大于特征数时，核形式反而比直接 Ridge 贵。两种视角谁划算，取决于 \(n\) 与 \(p\) 的对比。

\subsection{两个旋钮各控制什么}

核岭回归有两个主要超参数，作用机制完全不同。

\(\lambda\) 控制范数预算。正则项 \(\lambda n\,\alpha^{\trans}K\alpha\) 直接压在 \(\norm{f}_{\mathcal{H}_k}\) 上：\(\lambda\to0\) 时模型几乎只追求拟合，若 \(K\) 正定则退化为插值，训练误差为零；\(\lambda\to\infty\) 时 \(\widehat\alpha\to0\)，模型塌成常数。中间地带是偏差与方差的权衡，与``正则化从 Ridge 到 Lasso''中 \(\lambda\) 的角色一致，只是被收缩的不再是坐标方向的系数，而是 RKHS 几何下的``不光滑方向''——对应 \(K\) 的小特征值方向。

RBF 核的 \(\gamma\) 控制相似度的长度尺度。\(\gamma\) 大，\(k(x,z)\) 只在紧邻处非零，\(K\) 接近单位阵，几乎每个样本自立锚点，函数可以紧贴数据剧烈起伏；\(\gamma\) 小，远处样本也彼此相似，\(K\) 接近全一矩阵，允许的函数接近常数。两个旋钮联合决定有效容量：大 \(\gamma\) 配大 \(\lambda\) 是``精细但被强压''的模型，小 \(\gamma\) 配小 \(\lambda\) 是``本来就粗''的模型，验证误差的表面常常是二维的，只沿一条轴调参会漏掉谷底。

选择它们的纪律与``正则化从 Ridge 到 Lasso''强调的完全相同：预处理、核的选择与超参数的调节都必须包在每个训练折内部，用嵌套验证或从未参与选择的测试集评估整个流程。核化没有改变评估纪律，只改变了被评估的对象。

\subsection{代价：Gram 矩阵的三笔账}

核方法把特征维数从账单上抹掉了，代价是把样本量写上了账单：

\begin{itemize}
\tightlist
\item
  存储 \(K\)：\(O(n^{2})\)。\(n=10^{5}\) 时 \(K\) 有 \(10^{10}\) 个双精度数，约 80 GB，单台机器已放不下；
\item
  求解 \((K+\lambda nI)\alpha=y\)：稠密 Cholesky 分解 \(O(n^{3})\)。\(n=10^{5}\) 约需 \(10^{15}\) 次浮点运算；
\item
  预测：每个新样本要算 \(n\) 个核值，\(O(nd)\) 起步，且全部训练数据必须随模型保存。
\end{itemize}

二次存储与三次求解让朴素核方法在大数据上失效，这是它当年被更大规模的参数模型取代的直接原因之一。两条常用出路各用一句话概括：

\begin{itemize}
\tightlist
\item
  Nyström 近似：抽样 \(m\ll n\) 个锚点，只用 \(K\) 的对应列做低秩近似 \(K\approx CW^{-1}C^{\trans}\)，把存储与求解降到 \(O(nm)\) 与 \(O(m^{3})\) 量级，本质是承认 Gram 矩阵通常近似低秩；
\item
  随机特征（random features）：对平移不变核显式构造低维特征 \(z(x)\in\R^{m}\)，使 \(\ip{z(x)}{z(x')}\approx k(x,x')\)，然后回到原始形式训练线性模型——核被重新``显式化''，训练成本回到线性模型的水平。
\end{itemize}

两条路都在``核给的表达力''与``线性方法的规模''之间找付得起的折中，近似质量由 \(m\) 控制。

\subsection{为什么还要付这笔账}

如果只是为了拟合训练数据，核方法并非必需；真正为它定价的是泛化。``统计学习理论：训练误差对真风险有多乐观''一章的核心问题是训练误差比真风险乐观多少，答案取决于所用函数类的有效容量。对核方法，容量恰好由范数球

\[
\bigl\{\,f\in\mathcal{H}_k:\ \norm{f}_{\mathcal{H}_k}\le B\,\bigr\}
\]

刻画：这类函数集的复杂度量（如 Rademacher 复杂度）以 \(B\) 与核的谱衰减速度为界，\(\lambda\) 通过限制 \(B\) 直接收缩训练误差与真风险的差距。也就是说，正则项不是工程技巧，\(\norm{f}_{\mathcal{H}_k}\) 正是控制有效复杂度的那个范数——这与 SVM 中``最大化间隔等价于最小化 \(\norm{w}^{2}\)''是同一枚硬币的两面。

回看这一单元的主线：SVM 对偶里``数据只以内积出现''的观察，先给出核技巧这条替换规则，再要求一个严格的家（RKHS 与表示定理），最后以闭式解落地、以 \(O(n^{2})\) 与 \(O(n^{3})\) 标价。这个观察走得相当远——它统一了 SVM、Ridge 与 PCA 的核化形态，把``用范数控制容量''变成了可计算的对象；它走不到的地方，则标在大样本的账单上，以及``核本身需要人来选择''这个没有被定理消除的自由度里。
